%% notes-tex/figure-3-gate/sections/0-sketch.tex
%% SOURCE: figures/fig3-si-morphology-sketch.pdf is exported from
%% notes/assets/drawio/fig3-si-morphology-sketch.drawio.svg, written by
%% experiments/019-simb-multimodal/scripts/fig3_si_morphology_sketch_drawio.py (boxes, no data);
%% the counts in the panel boxes are read from results/morph_v24_dataset.json and
%% results/morph_v24_readout.json by that script's author and are repeated in
%% tables/morph_v24_dataset.tex and tables/morph_v24.tex (generated).

\section{The morphology SI figure, as a sketch to run experiments against\secstatus{todo}}
\label{sec:sketch}

\begin{keybox}
\textbf{Measured.} Panels a, b, c and f can be drawn now from committed result files
(\cref{tab:morph-strains,tab:morph-target,tab:morph-v24}, and the v20 readout).

\textbf{Not measured.} Panel d (a measured modality revealed, paired against the permuted
and the genotype-only controls) is in flight on \file{cabbi}; panel e (per feature) has no
dump behind it; panel g (the 1{,}011 isolates) is planned and likely not in time.
\end{keybox}

\tcfig{figures/fig3-si-morphology-sketch.pdf}{\textbf{The morphology SI figure, sketched
before its panels exist.} Each box names the quantity, the statistic and the sample the panel
will show, and the result file it will be drawn from; the status mark says whether that file
is committed. \textbf{a}, the three single-deletion panels and their overlaps (counts).
\textbf{b}, the target: per-feature replicate reliability, the 0.5 threshold that defines the
116 moving features, the ceiling. \textbf{c}, genotype only: the ceiling, the sequence
baselines and the cell graph transformer read off the perturbed CLS, twelve split seeds.
\textbf{d}, a measured modality revealed as input: the strain's own proteome (expression)
against another strain's and against genotype only on the same strains, paired by split
seed, beside the ridge from the same measurement. \textbf{e}, per feature: held-out Pearson
against reliability for the best revealed arm. \textbf{f}, the conditioning triangle: the
same contrast on the expression and proteome sides (v20). \textbf{g}, natural variation: the
1{,}011 isolates with Caudal 2024 expression and Muenzner 2024 proteome, planned.}{fig:si-sketch}

\pfstep{Every arm here is transductive in the sense the figure needs} The revealed modality
is a gene-level measurement of the \emph{same} strain, present on every training, validation
and test strain of the arm (\texttt{require\_modalities} keeps only genotypes carrying both
labels, \cref{tab:morph-strains}), and it enters through the observed-label channel on
every step with nothing masked, so the model is applied to a strain whose other modality
is known. Morphology itself can be a target and not an input: the observed-label channel
takes per-gene values, and CalMorph is one vector per cell population. The genotype-only
arms (\texttt{M\_cls}, \texttt{MP\_cls}) are the inductive controls on exactly the same
strains.

\pfstep{What the sketch asks of the runs} Panel d is the figure's claim and its only open
measurement: twelve paired differences per contrast, scored on the window rule, with the
ridge from the same measured modality as the bar (0.280 from the proteome, 0.324 from
expression, on the moving features). Panel c already says the genotype-only model sits at
0.046 $\pm$ 0.015 against a ceiling of 0.847, so the SI figure's honest reading is that
node-level measurements predict a cell-level phenotype and the genotype alone, through this
trunk, does not.

\pfstep{The 1{,}011 isolates} The same conditioning on natural variation needs the
Muenzner 2024 proteome built beside the served Caudal 2024 transcriptomes; morphology is
not measured on that panel, so the target there is the proteome from expression and the
reverse. It is first in the next-build band of the expansion document and is noted here so
the figure's slot exists; it is unlikely to be filled in this round.

%% GENERATED by experiments/019-simb-multimodal/scripts/morph_v24_dataset_table.py. Do not edit by hand.
%% SOURCE: the two stores' processed/phenotype_label_index.json and data_module_cache/index_seed_0.json;
%% results/morphology_noise_ceiling.json, morphology_feature_ceiling.csv, calmorph_moving_features.json,
%% proteome_morphology_covariation.json, expression_morphology_covariation.json, morphology_baselines_split_fig3_core/
\begin{table}[htbp]
  \centering
  \footnotesize
  \caption[The morphology round's strain sets]{The morphology round's strain sets. Genotype records per store and per label, the overlaps, the split seed 0 partition, and per arm the rows left after the training script's \texttt{require\_modalities} filter (the counts the runs print). A genotype is one deletion strain after deduplication and aggregation; a strain measured by two panels is one record carrying both labels. \src{experiments/019-simb-multimodal/scripts/morph_v24_dataset_table.py}}
  \label{tab:morph-strains}
  \begin{tabular}{@{}llrrrrrrr@{}}
    \toprule
    store & sources & genotypes & morphology & expression & proteome & morph.\ and expr.\ & morph.\ and prot.\ & all three \\
    \midrule
    \texttt{fig3\_core} & Kemmeren 2014, Sameith 2015, Ohya 2005, Costanzo 2016 and Kuzmin 2018/2020 fitness & 5,125 & 4,695 & 1,554 & -- & 1,439 & -- & -- \\
    \texttt{fig3\_morph\_proteome} & Kemmeren 2014, Sameith 2015, Messner 2023, Ohya 2005 & 4,977 & 4,695 & 1,554 & 4476 & 1,439 & 4281 & 1321 \\
    \midrule
    \multicolumn{9}{@{}l}{\emph{Per arm, split seed 0, after the label filter (train / validation / test rows)}}\\
    \multicolumn{2}{@{}l}{\texttt{M\_cls, M\_pool} on \texttt{fig3\_core}} & \multicolumn{7}{l}{morphology: 3,757 / 469 / 469} \\
    \multicolumn{2}{@{}l}{\texttt{CM\_expr, CM\_exprperm, CMP\_expr, CMP\_exprperm} on \texttt{fig3\_core}} & \multicolumn{7}{l}{expression, morphology: 1,161 / 132 / 146} \\
    \multicolumn{2}{@{}l}{\texttt{MP\_cls, CM\_prot, CM\_protperm} on \texttt{fig3\_morph\_proteome}} & \multicolumn{7}{l}{proteome, morphology: 3,464 / 416 / 401} \\
    \bottomrule
  \end{tabular}
\end{table}

\begin{table}[htbp]
  \centering
  \footnotesize
  \caption[The morphology target]{The morphology target and its references. CalMorph (Ohya 2005) gives one population-averaged profile per deletion and 122 wild-type replicates; a feature's reliability is one minus the replicate variance over the variance across mutants, its ceiling the square root, the best across-strain Pearson a perfect predictor of the signal can reach against one noisy measurement. The round predicts the moving features only. The ridge rows are the out-of-fold median Pearson per feature from the measured other modality on the strains both panels measured; the sequence rows are the best key selected on validation, four split seeds, all 278 features. \src{experiments/019-simb-multimodal/scripts/morph_v24_dataset_table.py}}
  \label{tab:morph-target}
  \begin{tabular}{@{}lr@{}}
    \toprule
    quantity & value \\
    \midrule
    deletion strains measured (Ohya 2005) & 4,718 \\
    wild-type replicates & 122 \\
    base CalMorph features & 281 \\
    after dropping the three degenerate & 278 \\
    moving features, reliability $\geq$ 0.5 (the target) & 116 \\
    target transform & Yeo-Johnson per feature, fit on the training rows \\
    ceiling, mean over all 278 features & 0.611 \\
    ceiling, mean over the moving features (minimum) & 0.847 (0.712) \\
    ridge from the measured proteome, moving features & 0.280 \\
    ridge from measured expression, moving features & 0.324 \\
    sequence bilinear ridge, all 278 features & 0.086 \\
    sequence neighbor average, all 278 features & 0.088 \\
    \bottomrule
  \end{tabular}
\end{table}

%% GENERATED by experiments/019-simb-multimodal/scripts/morph_v24_readout.py. Do not edit by hand.
%% SOURCE: W&B zhao-group/torchcell_019_morph_v24 run histories (IGB jobs 2427784 and 2427785);
%% results/morphology_feature_ceiling.csv; results/morphology_baselines_split_fig3_core/seed*.json
\begin{table}[htbp]
  \centering
  \footnotesize
  \caption[The morphology round v24]{The morphology round v24: the 116 moving CalMorph features read off the perturbed CLS, on fig3\_core, twelve split seeds (three for \texttt{M\_pool}). Window Pearson is the mean validation Pearson per feature over epochs 1{,}000 to 1{,}199, mean $\pm$ sd over the finished runs; the paired column is the per-split difference against the arm named, mean and number of splits above zero. \texttt{M\_cls} and \texttt{M\_pool} are validated on every calmorph strain, the \texttt{CM} arms on the strains carrying both labels, so only the two contrasts shown are paired. The replicate ceiling of the moving set is 0.847; the sequence baselines on all 278 features (four split seeds) are 0.086 (bilinear ridge) and 0.088 (neighbor average). \src{experiments/019-simb-multimodal/scripts/morph_v24_readout.py}}
  \label{tab:morph-v24}
  \begin{tabular}{@{}lp{40mm}rrlr@{}}
    \toprule
    arm & change & finished & collapsed & window Pearson & paired difference \\
    \midrule
    \texttt{M\_cls} & genotype only, CLS readout, all calmorph strains & 12/12 & 0 & 0.046 $\pm$ 0.015 & -- \\
    \texttt{M\_pool} & genotype only, CLS plus gene pool, all calmorph strains & 3/3 & 0 & 0.048 $\pm$ 0.025 & +0.015 (2/3) vs CLS alone \\
    \texttt{CM\_expr} & the strain's expression revealed, CLS readout (blind to it by construction) & 12/12 & 0 & 0.038 $\pm$ 0.025 & +0.000 (9/12) vs permuted \\
    \texttt{CM\_exprperm} & another strain's expression revealed, CLS readout (blind to it by construction) & 12/12 & 0 & 0.038 $\pm$ 0.026 & -- \\
    \texttt{CMP\_expr} & the strain's expression revealed, pooled readout, both-label strains & 12/12 & 0 & 0.122 $\pm$ 0.066 & +0.092 (12/12) vs permuted \\
    \texttt{CMP\_exprperm} & another strain's expression revealed, pooled readout, both-label strains & 12/12 & 0 & 0.030 $\pm$ 0.016 & -- \\
    \texttt{MP\_cls} & genotype only, pooled readout, proteome-and-morphology strains & 12/12 & 0 & 0.070 $\pm$ 0.029 & -- \\
    \texttt{CM\_prot} & the strain's proteome revealed, pooled readout, proteome-and-morphology strains & 3/3 & 0 & \textbf{0.133 $\pm$ 0.033} & +0.105 (1/1) vs permuted \\
    \texttt{CM\_protperm} & another strain's proteome revealed, pooled readout, proteome-and-morphology strains & 1/1 & 0 & 0.066 & -0.004 (0/1) vs genotype only \\
    \bottomrule
  \end{tabular}
\end{table}

